% Options for packages loaded elsewhere
% Options for packages loaded elsewhere
\PassOptionsToPackage{unicode}{hyperref}
\PassOptionsToPackage{hyphens}{url}
%
\documentclass[
  ignorenonframetext,
  aspectratio=169,
  russian,
]{beamer}
\newif\ifbibliography
\usepackage{pgfpages}
\setbeamertemplate{caption}[numbered]
\setbeamertemplate{caption label separator}{: }
\setbeamercolor{caption name}{fg=normal text.fg}
\beamertemplatenavigationsymbolshorizontal
% remove section numbering
\setbeamertemplate{part page}{
  \centering
  \begin{beamercolorbox}[sep=16pt,center]{part title}
    \usebeamerfont{part title}\insertpart\par
  \end{beamercolorbox}
}
\setbeamertemplate{section page}{
  \centering
  \begin{beamercolorbox}[sep=12pt,center]{section title}
    \usebeamerfont{section title}\insertsection\par
  \end{beamercolorbox}
}
\setbeamertemplate{subsection page}{
  \centering
  \begin{beamercolorbox}[sep=8pt,center]{subsection title}
    \usebeamerfont{subsection title}\insertsubsection\par
  \end{beamercolorbox}
}
% Prevent slide breaks in the middle of a paragraph
\widowpenalties 1 10000
\raggedbottom
\AtBeginPart{
  \frame{\partpage}
}
\AtBeginSection{
  \ifbibliography
  \else
    \frame{\sectionpage}
  \fi
}
\AtBeginSubsection{
  \frame{\subsectionpage}
}
\usepackage{iftex}
\ifPDFTeX
  \usepackage[T1]{fontenc}
  \usepackage[utf8]{inputenc}
  \usepackage{textcomp} % provide euro and other symbols
\else % if luatex or xetex
  \usepackage{unicode-math} % this also loads fontspec
  \defaultfontfeatures{Scale=MatchLowercase}
  \defaultfontfeatures[\rmfamily]{Ligatures=TeX,Scale=1}
\fi
\usepackage{lmodern}

\usetheme[]{default}
\ifPDFTeX\else
  % xetex/luatex font selection
\fi
% Use upquote if available, for straight quotes in verbatim environments
\IfFileExists{upquote.sty}{\usepackage{upquote}}{}
\IfFileExists{microtype.sty}{% use microtype if available
  \usepackage[]{microtype}
  \UseMicrotypeSet[protrusion]{basicmath} % disable protrusion for tt fonts
}{}


\usepackage{longtable,booktabs,array}
\newcounter{none} % for unnumbered tables
\usepackage{calc} % for calculating minipage widths
\usepackage{caption}
% Make caption package work with longtable
\makeatletter
\def\fnum@table{\tablename~\thetable}
\makeatother
\usepackage{graphicx}
\makeatletter
\newsavebox\pandoc@box
\newcommand*\pandocbounded[1]{% scales image to fit in text height/width
  \sbox\pandoc@box{#1}%
  \Gscale@div\@tempa{\textheight}{\dimexpr\ht\pandoc@box+\dp\pandoc@box\relax}%
  \Gscale@div\@tempb{\linewidth}{\wd\pandoc@box}%
  \ifdim\@tempb\p@<\@tempa\p@\let\@tempa\@tempb\fi% select the smaller of both
  \ifdim\@tempa\p@<\p@\scalebox{\@tempa}{\usebox\pandoc@box}%
  \else\usebox{\pandoc@box}%
  \fi%
}
% Set default figure placement to htbp
\def\fps@figure{htbp}
\makeatother



\ifLuaTeX
\usepackage[bidi=basic,shorthands=off,provide=*]{babel}
\else
\usepackage[bidi=default,shorthands=off,provide=*]{babel}
\fi
\ifLuaTeX
  \usepackage{selnolig} % disable illegal ligatures
\fi


\setlength{\emergencystretch}{3em} % prevent overfull lines

\providecommand{\tightlist}{%
  \setlength{\itemsep}{0pt}\setlength{\parskip}{0pt}}



 

\usepackage[]{csquotes}

\IfFileExists{plex-otf.sty}{
  %% Full TeXlive
  % \usepackage[%
  %   % math,
  %   RM={Scale=0.94},SS={Scale=0.94},SScon={Scale=0.94},TT={Scale=MatchLowercase,FakeStretch=0.9},DefaultFeatures={Ligatures=Common}
  % ]{plex-otf}
}{
  %% TinyTeX
  \usepackage{libertine}
}

%%% Load theme
% https://deic.uab.cat/~iblanes/beamer_gallery/
\IfFileExists{beamerthemegotham.sty}{
  %% Full TeXlive
  \usetheme{gotham}
  \gothamset{
    numbering=totalpagenumber,
    parttocframe default=off,
    sectiontocframe default=off,
    subsectiontocframe default=off,
  }
}{
  %% TinyTeX
  \usetheme{Madrid}
}
\metroset{progressbar=frametitle,sectionpage=none,numbering=fraction}
\usepackage{fontspec}
\usepackage{polyglossia}
\setdefaultlanguage{russian}
\setotherlanguage{english}
\defaultfontfeatures{Ligatures=TeX}
\setsansfont{DejaVu Sans}
\setmainfont{DejaVu Serif}
\setmonofont{DejaVu Sans Mono}
\makeatletter
\@ifpackageloaded{caption}{}{\usepackage{caption}}
\AtBeginDocument{%
\ifdefined\contentsname
  \renewcommand*\contentsname{Содержание}
\else
  \newcommand\contentsname{Содержание}
\fi
\ifdefined\listfigurename
  \renewcommand*\listfigurename{Список иллюстраций}
\else
  \newcommand\listfigurename{Список иллюстраций}
\fi
\ifdefined\listtablename
  \renewcommand*\listtablename{Список таблиц}
\else
  \newcommand\listtablename{Список таблиц}
\fi
\ifdefined\figurename
  \renewcommand*\figurename{Рисунок}
\else
  \newcommand\figurename{Рисунок}
\fi
\ifdefined\tablename
  \renewcommand*\tablename{Таблица}
\else
  \newcommand\tablename{Таблица}
\fi
}
\@ifpackageloaded{float}{}{\usepackage{float}}
\floatstyle{ruled}
\@ifundefined{c@chapter}{\newfloat{codelisting}{h}{lop}}{\newfloat{codelisting}{h}{lop}[chapter]}
\floatname{codelisting}{Список}
\newcommand*\listoflistings{\listof{codelisting}{Листинги}}
\makeatother
\makeatletter
\makeatother
\makeatletter
\@ifpackageloaded{caption}{}{\usepackage{caption}}
\@ifpackageloaded{subcaption}{}{\usepackage{subcaption}}
\makeatother

\usepackage{bookmark}
\IfFileExists{xurl.sty}{\usepackage{xurl}}{} % add URL line breaks if available
\urlstyle{same}
\hypersetup{
  pdftitle={Лабораторная работа №1},
  pdfauthor={Лопатченко Полина Андреевна},
  pdflang={ru-RU},
  hidelinks,
  pdfcreator={LaTeX via pandoc}}


\title{\texorpdfstring{Лабораторная работа №1}{Лабораторная работа №1}}
\subtitle{\texorpdfstring{Установка ОС Linux (Fedora Sway) в
VirtualBox}{Установка ОС Linux (Fedora Sway) в VirtualBox}}
\author{Лопатченко Полина Андреевна}
\date{2026-03-04}
\institute{Российский университет дружбы народов, Москва, Россия}

\begin{document}
\frame{\titlepage}


\begin{frame}{Докладчик}
\protect\phantomsection\label{ux434ux43eux43aux43bux430ux434ux447ux438ux43a}
\begin{columns}[c]
\begin{column}{0.7\linewidth}
\begin{itemize}[<+->]
\tightlist
\item
  Лопатченко Полина Андреевна
\item
  Студент НКАбд-04-25
\item
  Российский университет дружбы народов
\item
  \href{mailto:1032253529@rudn.ru}{\nolinkurl{1032253529@rudn.ru}}
\end{itemize}
\end{column}

\begin{column}{0.3\linewidth}
\end{column}
\end{columns}

\#Информация
\end{frame}

\begin{frame}[fragile]{Цель и задачи}
\protect\phantomsection\label{ux446ux435ux43bux44c-ux438-ux437ux430ux434ux430ux447ux438}
\textbf{Цель:} установить Fedora Sway в VirtualBox и выполнить базовую
настройку системы.

\textbf{Задачи:}

\begin{itemize}[<+->]
\tightlist
\item
  создать виртуальную машину и установить ОС;
\item
  выполнить первичную настройку (обновления, инструменты, SELinux,
  раскладка, hostname);
\item
  установить ПО для подготовки отчётов (\texttt{pandoc}, TeX Live);
\item
  получить системную информацию через \texttt{dmesg}.
\end{itemize}
\end{frame}

\section{Ход
работы}\label{ux445ux43eux434-ux440ux430ux431ux43eux442ux44b}

\begin{frame}{1) Создание виртуальной машины}
\protect\phantomsection\label{ux441ux43eux437ux434ux430ux43dux438ux435-ux432ux438ux440ux442ux443ux430ux43bux44cux43dux43eux439-ux43cux430ux448ux438ux43dux44b}
Создала новую виртуальную машину в графическом интерфейсе и задал имя ВМ
по логину.

\begin{figure}

\centering{

\includegraphics[width=0.3\linewidth,height=\textheight,keepaspectratio]{image/1.png}

}

\caption{\label{fig-001}Создание новой виртуальной машины}

\end{figure}%
\end{frame}

\begin{frame}{2) Настройка ресурсов}
\protect\phantomsection\label{ux43dux430ux441ux442ux440ux43eux439ux43aux430-ux440ux435ux441ux443ux440ux441ux43eux432}
Указала объём оперативной памяти и размер диска 80 ГБ.

\begin{center}
\begin{minipage}{0.30\linewidth}
\centering
\includegraphics[width=\linewidth]{image/2.png}\\
\footnotesize Память виртуальной машины
\end{minipage}\hfill
\begin{minipage}{0.30\linewidth}
\centering
\includegraphics[width=\linewidth]{image/4.png}\\
\footnotesize Размер диска 80 ГБ
\end{minipage}
\end{center}
\end{frame}

\begin{frame}{3) Подключение установочного ISO}
\protect\phantomsection\label{ux43fux43eux434ux43aux43bux44eux447ux435ux43dux438ux435-ux443ux441ux442ux430ux43dux43eux432ux43eux447ux43dux43eux433ux43e-iso}
Подключила ISO-образ Fedora Sway.

\begin{figure}

\centering{

\includegraphics[width=0.4\linewidth,height=\textheight,keepaspectratio]{image/2.png}

}

\caption{\label{fig-002}Подключение ISO Fedora Sway}

\end{figure}%
\end{frame}

\begin{frame}{4) Параметры загрузки и виртуализации}
\protect\phantomsection\label{ux43fux430ux440ux430ux43cux435ux442ux440ux44b-ux437ux430ux433ux440ux443ux437ux43aux438-ux438-ux432ux438ux440ux442ux443ux430ux43bux438ux437ux430ux446ux438ux438}
Включила ускорение 3D/контроллер VMSVGA и поддержку UEFI. \#\# 5)
Установка Fedora Sway Установила систему на диск и завершил установку.

\begin{figure}

\centering{

\includegraphics[width=0.78\linewidth,height=\textheight,keepaspectratio]{image/2.png}

}

\caption{\label{fig-002}Завершение установки Fedora Sway}

\end{figure}%
\end{frame}

\section{Первичная настройка
Fedora}\label{ux43fux435ux440ux432ux438ux447ux43dux430ux44f-ux43dux430ux441ux442ux440ux43eux439ux43aux430-fedora}

\begin{frame}[fragile]{Обновление и инструменты}
\protect\phantomsection\label{ux43eux431ux43dux43eux432ux43bux435ux43dux438ux435-ux438-ux438ux43dux441ux442ux440ux443ux43cux435ux43dux442ux44b}
Установила \texttt{development-tools} и обновила пакеты.

\begin{figure}

\centering{

\includegraphics[width=0.6\linewidth,height=\textheight,keepaspectratio]{image/3.png}

}

\caption{\label{fig-003}Обновление пакетов}

\end{figure}%
\end{frame}

\begin{frame}[fragile]{SELinux и hostname}
\protect\phantomsection\label{selinux-ux438-hostname}
Перевёла SELinux в режим \texttt{permissive} и установил имя хоста по
требованиям именования.

\begin{figure}

\centering{

\includegraphics[width=0.3\linewidth,height=\textheight,keepaspectratio]{image/5.png}

}

\caption{\label{fig-005}SELinux permissive}

\end{figure}%

\begin{figure}

\centering{

\includegraphics[width=0.3\linewidth,height=\textheight,keepaspectratio]{image/6.png}

}

\caption{\label{fig-006}Установка hostname}

\end{figure}%
\end{frame}

\begin{frame}{Раскладка клавиатуры (Sway + system)}
\protect\phantomsection\label{ux440ux430ux441ux43aux43bux430ux434ux43aux430-ux43aux43bux430ux432ux438ux430ux442ux443ux440ux44b-sway-system}
Создала пользовательский конфиг Sway и отредактировал системный конфиг
клавиатуры.

\begin{figure}

\centering{

\includegraphics[width=0.6\linewidth,height=\textheight,keepaspectratio]{image/4.png}

}

\caption{\label{fig-004}Системный конфиг клавиатуры}

\end{figure}%
\end{frame}

\section{ПО для отчётов и обмен
файлами}\label{ux43fux43e-ux434ux43bux44f-ux43eux442ux447ux451ux442ux43eux432-ux438-ux43eux431ux43cux435ux43d-ux444ux430ux439ux43bux430ux43cux438}

\begin{frame}[fragile]{Pandoc и TeX Live}
\protect\phantomsection\label{pandoc-ux438-tex-live}
Установила \texttt{pandoc} и TeX Live для сборки PDF.

\begin{figure}

\centering{

\includegraphics[width=0.4\linewidth,height=\textheight,keepaspectratio]{image/7.png}

}

\caption{\label{fig-007}Установка pandoc}

\end{figure}%

\begin{figure}

\centering{

\includegraphics[width=0.4\linewidth,height=\textheight,keepaspectratio]{image/8.png}

}

\caption{\label{fig-008}Установка TeX Live}

\end{figure}%
\end{frame}

\section{Домашнее задание
(dmesg)}\label{ux434ux43eux43cux430ux448ux43dux435ux435-ux437ux430ux434ux430ux43dux438ux435-dmesg}

\begin{frame}{Ключевая системная информация (1/3)}
\protect\phantomsection\label{ux43aux43bux44eux447ux435ux432ux430ux44f-ux441ux438ux441ux442ux435ux43cux43dux430ux44f-ux438ux43dux444ux43eux440ux43cux430ux446ux438ux44f-13}
Определила версию ядра Linux. Определила частоту и модель процессора.

\begin{figure}

\centering{

\includegraphics[width=0.92\linewidth,height=\textheight,keepaspectratio]{image/9.png}

}

\caption{\label{fig-009}Версия ядра Linux}

\end{figure}%

\begin{figure}

\centering{

\includegraphics[width=0.92\linewidth,height=\textheight,keepaspectratio]{image/9.png}

}

\caption{\label{fig-009}Частота процессора}

\end{figure}%

\begin{figure}

\centering{

\includegraphics[width=0.92\linewidth,height=\textheight,keepaspectratio]{image/9.png}

}

\caption{\label{fig-009}Модель процессора}

\end{figure}%
\end{frame}

\begin{frame}{Ключевая системная информация (2/3)}
\protect\phantomsection\label{ux43aux43bux44eux447ux435ux432ux430ux44f-ux441ux438ux441ux442ux435ux43cux43dux430ux44f-ux438ux43dux444ux43eux440ux43cux430ux446ux438ux44f-23}
Определила объём доступной памяти и тип гипервизора.

\begin{figure}

\centering{

\includegraphics[width=0.92\linewidth,height=\textheight,keepaspectratio]{image/10.png}

}

\caption{\label{fig-010}Доступная оперативная память}

\end{figure}%

\begin{figure}

\centering{

\includegraphics[width=0.92\linewidth,height=\textheight,keepaspectratio]{image/11.png}

}

\caption{\label{fig-011}Тип гипервизора}

\end{figure}%
\end{frame}

\begin{frame}{Ключевая системная информация (3/3)}
\protect\phantomsection\label{ux43aux43bux44eux447ux435ux432ux430ux44f-ux441ux438ux441ux442ux435ux43cux43dux430ux44f-ux438ux43dux444ux43eux440ux43cux430ux446ux438ux44f-33}
Определил тип файловой системы корневого раздела и последовательность
монтирования ФС.

\begin{center}
\begin{minipage}{0.48\linewidth}
\centering
\includegraphics[width=\linewidth]{image/11.png}\\
\footnotesize Тип ФС корневого раздела
\end{minipage}\hfill
\begin{minipage}{0.48\linewidth}
\centering
\includegraphics[width=\linewidth]{image/10.png}\\
\footnotesize Последовательность монтирования ФС
\end{minipage}
\end{center}
\end{frame}

\section{Итоги}\label{ux438ux442ux43eux433ux438}

\begin{frame}[fragile]{Выводы}
\protect\phantomsection\label{ux432ux44bux432ux43eux434ux44b}
\begin{itemize}[<+->]
\tightlist
\item
  Создана ВМ в VirtualBox и установлена Fedora Sway.
\item
  Выполнена первичная настройка: обновления, инструменты, SELinux,
  раскладка, hostname.
\item
  Установлены \texttt{pandoc} и TeX Live; настроена общая папка для
  обмена файлами.
\item
  По \texttt{dmesg} получены параметры загрузки и сведения об окружении
  (CPU, память, гипервизор, ФС). :::
\end{itemize}
\end{frame}




\end{document}
